%
\againframe{plan}
%
\begin{frame}{Panchabhuta -- Five elements}
	\begin{itemize}
		\item Etymology is Pancha (Sanskrit for five) and bhuta (Sanskrit for elements) 
		\item Fundamental idea in Indian philosophy, Ayurveda, Yoga, and traditional sciences. It explains the composition of the universe, including the human body and mind, in terms of five basic elements.
		\item The Panchabhuta are considered the building blocks of the cosmos. Everything in the universe (macrocosm) and in the human body (microcosm) is made of these five elements in varying proportions.
	\end{itemize}
	%
\end{frame}
%
\begin{frame}{Panchabhuta}
\begin{footnotesize}
		\begin{table}
		\begin{tabular}{l l l l}
			\toprule
			Bhuta	&	Translation 	&	Associated Sense	&	Body part	\\
			\midrule
			Ākāśa & 	Space/Ether		&	Hearing (śabda)	&
			Cavities (Mouth)	\\
			&&& Channels (Nostrils)	\\
			Vāyu	&	Air	&	Touch (sparśa)	&	Spaces (Thorax, Abdomen)	\\
			Agni/Tejas	&	Fire	&	Sight (rūpa)	& digestion, vision \\
			&&& body temperature  \\
			&&& intelligence	\\
			Āpa / Jala	&	Water	&	Taste (rasa)	&	blood, saliva, digestive juices	\\
			&&& plasma, reproductive fluids	\\
			Pṛthvī 	&	Earth	& 	Smell (gandha)	&	bones, teeth, flesh, \\
			&&&	muscles, skin, nails, hair	\\
			\bottomrule
		\end{tabular}
		\caption{{\tiny 1. Ākāśa represents emptiness, vastness, and the medium that allows sound to travel. 2. Vāyu symbolizes movement, motion, and dynamism.
		3. Agni represents transformation, heat, energy, and metabolism.
		4. Jala symbolizes fluidity, cohesion, and life-sustaining properties.
		5. Pṛthvī represents solidity, stability, and structure.}}
	\end{table}
\end{footnotesize}
\end{frame}
%
\begin{frame}{}
	\begin{huge}
		\begin{center}
			End of Unit 1
		\end{center}
	\end{huge}

\end{frame}